\documentclass{article}

\usepackage{tabularx}
\usepackage{booktabs}
\usepackage{amsmath}
\usepackage{hyperref}
\hypersetup{colorlinks=true, linkcolor=blue, urlcolor=blue}

\title{Reflection and Traceability Report on \progname}

\author{\authname}

\date{}

\input{../Comments.text}
\input{../Common.text}

\begin{document}

\maketitle

% -----------------------------------------------------------------------
\section{Changes in Response to Feedback}
% -----------------------------------------------------------------------

\subsection{SRS and Hazard Analysis}

The SRS went through two rounds of revision based on feedback from a
peer reviewer and Dr.\ Smith.
The main changes were: adding a complete Introduction section,
clarifying the intended user background, making modelling and scope
decisions explicit, expanding the traceability matrices, adding software
constraints, and fixing notation inconsistencies.

The most significant self-initiated change was restricting field regions
to a uniform grid.
The original free-form corner specification was replaced by an
equal-size, grid-aligned layout.
This simplified both the implementation and verification.
The charge constraint was also broadened from $q>0$ to
$|q|\le q_\text{max}$ to support negative charges.

The Hazard Analysis received no feedback requiring major changes.

\subsection{Design and Design Documentation}

The Trajecto implementation is generated by Drasil.
The design documents describe the generated module structure.

\subsection{VnV Plan and Report}

The VnV Plan went through two rounds of revision based on feedback from
a peer reviewer and Dr.\ Smith.
Key changes were: adding new functional test cases, adding a dedicated
input-validation section, adding a coverage table, replacing speculative
error bounds with actual test results, and correcting the Software
Validation statement.

% -----------------------------------------------------------------------
\section{Extras}
% -----------------------------------------------------------------------

No extras were planned for this project.

% -----------------------------------------------------------------------
\section{Design Iteration (LO11)}
% -----------------------------------------------------------------------

The most significant design iteration was in how field regions are
specified.
The original SRS allowed freely specified corner coordinates.
This created ambiguity around overlapping regions and gaps.
After review during implementation, regions were restricted to a uniform
grid.
This made the input well-defined and enabled analytic verification.

A second iteration involved the charge constraint.
The original $q>0$ requirement was updated to $|q|\le q_\text{max}$.
Negative charges are physically meaningful and are already handled
correctly by the sign of $q$ in the Lorentz force.

% -----------------------------------------------------------------------
\section{Design Decisions (LO12)}
% -----------------------------------------------------------------------

\textbf{Limitations.}
The implementation is generated by Drasil.
Drasil currently only supports Dormand--Prince RK45 as the ODE solver.
This was not user-configurable, but it is well-suited for smooth
trajectory problems.

\textbf{Assumptions.}
Motion is restricted to the $xy$-plane to keep the problem tractable.
The uniform-grid region assumption was made to enable analytic
verification and simplify input parsing.

\textbf{Constraints.}
Physical constant bounds were added as software constraints.
They prevent the solver from being called with numerically degenerate
values.
Most are enforced as soft warnings rather than hard errors, to
accommodate research users who may explore edge cases.

% -----------------------------------------------------------------------
\section{Economic Considerations (LO23)}
% -----------------------------------------------------------------------

Trajecto is a research and educational tool.
There is no commercial market for a bespoke charged-particle simulator
at this scale.

Trajecto's value is as a demonstration case for the Drasil
knowledge-capture framework.
Potential users are researchers working in scientific software quality
and charged-particle physics.

% -----------------------------------------------------------------------
\section{Reflection on Project Management (LO24)}
% -----------------------------------------------------------------------

\subsection{How Does Your Project Management Compare to Your Development Plan}

The workflow plan of using GitHub Issues to track feedback was followed
consistently.
The technology stack (LaTeX, Python, Drasil, GitHub) matched the
Development Plan.

\subsection{What Went Well?}

Using GitHub Issues to log each feedback item worked well.
Every piece of feedback was recorded and each fix was tied to a commit.
This made writing the traceability sections straightforward.

Choosing the Drasil-generated Python implementation early saved
significant debugging time.

\subsection{What Went Wrong?}

Some VnV test cases were written based on SRS assumptions that the
generated code did not fully satisfy.
Several input-validation tests could not be executed.
The grid-based input format cannot express the invalid region topologies
those tests assumed.
This only became apparent during test execution.
Two input constraints were also absent from the generated code.
They were discovered as gaps during testing rather than during design.

\subsection{What Would You Do Differently Next Time?}

When writing the VnV Plan, cross-check each test case against the
Drasil specification as well as the SRS.
The specification determines what the generated code will enforce, not
just the SRS.
Running a quick prototype of the input parser earlier would also help.
It would reveal which constraint checks are reachable before the VnV
Plan is finalised.

% -----------------------------------------------------------------------
\section{Ethics and Equity (LO18)}
% -----------------------------------------------------------------------

Trajecto simulates charged-particle physics.
It has no human-facing output and handles no sensitive data.
There are no specific equity concerns.

The input is plain text and error messages are in plain English.
The output is a human-readable file.
Any user with a terminal can run the tool without special
hardware or assistive technology.

% -----------------------------------------------------------------------
\section{Reflection on Capstone}
% -----------------------------------------------------------------------

\subsection{Which Courses Were Relevant}

\begin{itemize}
  \item \textbf{Physics} ---
        the Lorentz force law and equations of motion are central to
        the project.
  \item \textbf{Numerical Methods} ---
        understanding RK45 accuracy was needed to write meaningful
        pseudo-oracle tests.
  \item \textbf{Software Engineering (CAS~741)} ---
        the overall framework: SRS, VnV Plan, hazard analysis, and
        traceability.
  \item \textbf{Python programming} ---
        implementing and debugging the test scripts.
\end{itemize}

\subsection{Knowledge/Skills Outside of Courses}

\begin{itemize}
  \item \textbf{Drasil framework} ---
        understanding the knowledge system and the structure of
        its generated code.
  \item \textbf{LaTeX} ---
        writing large structured documents with multiple references and
        files.
  \item \textbf{GitHub Issues as a traceability tool} ---
        using issue tables linked to commits as a systematic record of
        changes based on feedback.
\end{itemize}

\end{document}
